% Revision notes — one topic worked over its whole spread of sources, in a topic folder under
% 30 Revision/. The lecture pattern with more sources, and questions reached by pointer.
% Copy it there and replace every <...> placeholder; the input line below needs no edit.
\documentclass[a4paper,10pt]{article}
\IfFileExists{preamble.tex}{\input{preamble.tex}}{\input{../../templates/preamble.tex}}
\begin{document}
\ArtifactHeader{\ModuleCode}{<Source-map key>}{<the topic, from the Source map's topics>}%
  {<YYYY-MM-DD>}{<every unit, sheet and paper this topic spans, by pointer>}

\section{What this topic is}
<The topic in the shape it will be examined in, not the order it was taught in. Revision earns its
place by re-cutting the material, so a walkthrough copied here has done nothing.>

\section{What has to be known cold}

\begin{defn}[<term>]
<The definitions the rest of the topic is stated in terms of.>
\end{defn}

\begin{thm}[<name>]
<The result the topic turns on, hypotheses included.>
\end{thm}

\begin{prop}
<A smaller result that gets used constantly and is never named in the lectures.>
\end{prop}

\begin{remark}
<How the units that taught this fit together — usually the thing no single unit could say.>
\end{remark}

\begin{warning}
<The confusion this topic reliably produces, taken from the records and the Revisit register.>
\end{warning}

\section{Where it is examined}
\begin{itemize}
  \item <A sheet question, by pointer — never copied here.>
  \item <A past-paper question, by pointer, with what makes it the representative one.>
\end{itemize}

\begin{checkyourself}
  \item <The derivation to be able to produce from nothing.>
  \item <The question that would show this topic is closed.>
\end{checkyourself}
\end{document}
